\scanpage{100}

Conditions nécessaires pour qu'un $\lambda\in\bar Q$ soit valeur propre d'un \unsure{frobénius}, pour une variété projective lisse définie sur le corps $\mathbb F_q$, $q=p^n$:

\deleted{1) $\lambda\in Q'(i)$

2) \deleted{\unsure{$\lambda$}} $\lambda^2$ est de la forme $q^iu$, où $u\in Q'(i)$ est une unité exc. $p$.

Si $\mathfrak p,\mathfrak p'$ sont deux idéaux premiers de}

1) \fbox{$\lambda$ entier algébrique}

2) \fbox{$\lambda\in Q'(\sqrt{-1})$}

3) $\lambda^{\unsure{2}}=q^{i/2}u$, \inserted{$i\in Z$,} avec $u$ une unité \inserted{absolue} exc. $p$ (cela détermine évidemment de façon unique $i$ et $u$)

\inserted{\fbox{$\exists i$ tel que $|\lambda_\alpha|=q^{i/2}$ pour \unsure{tt|$\forall$} conjugué $\lambda_\alpha$} !}

\margin{N.B. $q^{i/2}\in Q'$!}

\margin{\uline{N.B.} Les conditions 2) 3) 4) sont ``\unsure{linéaires}'', la condition 1) est seulement \unsure{additive} ... On désigne par $M(q^{\unsure{n}})$ le groupe de $Q'(i)^*$ défini par \deleted{\unsure{2)}} 3) 4).}

\deleted{4) Si $\mathfrak p,\mathfrak p'$ sont deux idéaux premiers de $Q'(\sqrt{-1})$ au dessus du même idéal premier de $Q'$, on a
\figure{\begin{tikzpicture}[x=1cm,y=1cm]
\node at (0,1.5) {$v_{\mathfrak p}(u)+v_{\mathfrak p'}(u)=0$};
\node at (3.4,1.5) {i.e.};
\node[draw,rectangle] (f) at (7,1.5) {$v_{\mathfrak p}(\lambda)+v_{\mathfrak p'}(\unsure{\lambda^*|\lambda})=\dfrac{in}{2}$};
\node[align=center] at (7,0) {$v_{\mathfrak p}$ normalisés tels que\\$v_{\mathfrak p}(p)=1$.};
\node[draw,ellipse,align=center] (c) at (2,-.4) {Conséquence\\de \uline{2 et 3}\\!!!};
\draw[->,double] (c.north east) -- (f.south west);
\end{tikzpicture}}}

L'entier $i$ qui intervient ici est appelé le \uline{poids} de $\lambda$ (relativement à $q=p^n$). \deleted{\unsure{Le fait que le poids est}} Si les conditions \deleted{\unsure{est donc}} analogues sont satisfaites pour $q'=p^m$, elles le sont pour $q$, et le poids $i$ rel. à $q$ est égal à $mi'$, où $i'$ est le poids par rapport à $q'$.

\margin{\unsure{Th. d'existence en $\mathrm{Fr}$ sur $\mathbb F_q$}}

\margin{\figure{\begin{tikzpicture}\draw (0,0)--(0,1.1);\draw (.06,0)--(.06,1.1);\draw (1.35,0)--(1.35,1.15);\end{tikzpicture}}}

\uline{Conjecture} \deleted{\unsure{est elle}} \uline{forte} \uline{Les conditions nécessaires précédentes sont suffisantes}

Cela déterminerait donc explicitement, en termes arithmétiques simples, le groupe de Galois motivique \inserted{$G$} sur $\mathbb F_q$. Vérifions que cela donne bien

\mspage{99}
